% Generated draft: requires mathematical and source review.
\documentclass[11pt]{article}
\usepackage[margin=25mm]{geometry}
\usepackage{fontspec}
\setmainfont{Latin Modern Roman}
\usepackage{amsmath,amssymb,mathtools,unicode-math}
\setmathfont{Latin Modern Math}
\usepackage{enumitem,needspace,xurl,hyperref}
\setlength{\parindent}{0pt}
\setlength{\parskip}{5pt}
\setlength{\emergencystretch}{3em}
\newcommand{\N}{\mathbb{N}}
\newcommand{\Z}{\mathbb{Z}}
\newcommand{\Q}{\mathbb{Q}}
\newcommand{\R}{\mathbb{R}}
\newcommand{\C}{\mathbb{C}}
\newcommand{\F}{\mathbb{F}}
\newcommand{\A}{\mathbb{A}}
\newcommand{\PP}{\mathbb{P}}
\newcommand{\E}{\mathbb{E}}
\newcommand{\eps}{\varepsilon}
\newcommand{\dd}{\,\mathrm{d}}
\begin{document}
\section*{UnsolvedMath ID: 20002866}
\textbf{Source problem code:} catalogue entry
\subsection{Definitions and mathematical statement}
Let \(\mathbb B^n=\{z\in\C^n:\|z\|<1\}\), and let \(f:\mathbb B^n\to\mathbb B^N\) be a proper rational holomorphic map. Properness means compact inverse images of compact sets. Write \(f=P/q\), with a common polynomial denominator, the components of \(P\) and \(q\) having no common nonconstant factor, and set
\[
\deg f=\max(\deg P_1,\ldots,\deg P_N,\deg q)=d.
\]
The source question, attributed to J. D'Angelo, is whether there exists a polynomial proper map \(p:\mathbb B^n\to\mathbb B^N\) of degree exactly \(d\), and a continuous family \((f_t)_{0\le t\le1}\) of proper holomorphic maps into the \emph{same} target ball, such that
\[
f_0=f,\qquad f_1=p.
\]
For definiteness, continuity may be read in the compact-open topology on holomorphic maps, equivalently locally uniform continuity in the parameter. The original wording specifies properness of every \(f_t\), but does not specify a topology or explicitly require rational intermediate maps. It requires equal degrees of the endpoints; a path with \(\deg f_t=d\) for every \(t\) is a stronger conclusion.

The supplied research aid reports fixed-target rational homotopies of constant algebraic degree for maps spherically equivalent to a polynomial, where spherical equivalence means composition with automorphisms of the source and target balls. It reports in particular that degree-two maps from \(\mathbb B^2\) satisfy the requested conclusion. It further describes an automorphism-slot denominator-removal lemma: an automorphism appearing in a factor multiplied by a polynomial \(h\) can be contracted to a unitary map while preserving properness; constant degree requires the stated degree dominance and nondivisibility of \(h\) by the affine denominator. The precise slot construction is not supplied here and cannot be taken as a lemma for arbitrary denominators.

As a further example the aid reports a constant-degree path \(F_{(1-t)a}\) in the standard degree-three Faran--Huang--Ji--Zhang family, which is not spherically equivalent to a polynomial, ending at
\[
p(z,w)=(z^2,\sqrt2\,zw,zw^2,w^3)\in\C^4.
\]
No formula or admissible parameter range for \(F_a\) is supplied in the record. Finally the aid reports a separate stabilized result: after including \(f\) by zero coordinates in a ball of target dimension
\[
M=n+\max\{N,d(n-1)+1\},
\]
a proper rational homotopy reaches a degree-\(d\) Whitney polynomial, with intermediate degrees at most \(d+1\). This statement permits a larger target and therefore does not answer the original fixed-target problem. The trailing source words ``Squared norms'' are the next heading, not a further hypothesis.
\subsection{Short English statement}
Can every rational proper ball map be deformed through proper maps in its fixed target to a polynomial of the same degree? The supplied aid gives special constant-degree homotopies and a result after increasing the target dimension; the stabilized result leaves the original requirement unresolved.
\subsection{Sources}
\begin{itemize}
\item J. D'Angelo, Problem 5, \emph{Mappings between balls}, 2010 AIM workshop \emph{Emerging applications of complexity for CR mappings}, \url{https://aimath.org/WWN/crmappings/crmappings.pdf}; attribution, year and section are supplied by the record verification.
\item The research summary names the Faran--Huang--Ji--Zhang degree-three family but gives no full bibliographic reference or formula for that family.
\item Local UnsolvedMath registry, permanent ID 20002866, exact code \texttt{catalogue entry}; canonical location \texttt{aim-several-complex-variables-notes.json}, notes[23], dated 14 August 2026; statement labeled exact.
\item Research attempt 1: partial result, problem partially solved at the run, rigor label proved, novelty confidence low; registry status partially solved. The homotopy constructions are reported by a machine-generated aid, not peer reviewed, and require independent verification before citation.
\end{itemize}
\end{document}
